% Source: https://artofproblemsolving.com/community/c3h2516562_weird_desmos_graph
% Stable identity: aops-topic-2516562-q01
\textbf{(a)} For $x,y>0$, taking logarithms,
\[x^y=y^x\iff y\ln x=x\ln y\iff g(x)=g(y),\qquad g(t)=\dfrac{\ln t}{t}.\]
Since $g'(t)=\dfrac{1-\ln t}{t^2}$, $g$ is strictly increasing on $(0,e]$ and strictly decreasing on $[e,\infty)$. So $g(x)=g(y)$ with $x<y$ is impossible if both lie in $(0,e]$ or both in $[e,\infty)$; hence $x<e<y$. Then $g(y)>0$ (as $y>1$), so $g(x)>0$, i.e. $x>1$.

\textbf{(b)} With $y=kx$: $x^{kx}=(kx)^x\iff x^k=kx\iff x^{k-1}=k$ (taking $x$-th roots, as $x>0$). So
\[x=k^{\dfrac{1}{k-1}},\qquad y=kx=k^{\dfrac{k}{k-1}}.\]
For example $k=2$ gives $(2,4)$: $2^4=4^2$.

\textbf{(c)} By (a), points of the second curve with $x<y$ satisfy $x<e<y$, and by symmetry the points with $x>y$ satisfy $y<e<x$; so the only point of the line $y=x$ that the curve can approach is $(e,e)$. It does approach it: as $k\to1^+$, put $k=1+h$ with $h\to0^+$; then
\[x=(1+h)^{\dfrac1h}\to e,\qquad y=kx\to e.\]
At $(e,e)$ the equation holds ($e^e=e^e$), and the two branches cross there. So the curve meets the line at $\boxed{(e,e)}$.

\begin{asy}
usepackage("amsmath");
size(5.5cm);
defaultpen(fontsize(10pt));
import graph;
real L=6.2;
draw((0,0)--(L,0),Arrow(5)); draw((0,0)--(0,L),Arrow(5));
label("$x$",(L,0),E); label("$y$",(0,L),N);
for(int i=1;i<=6;++i){draw((i,-0.06)--(i,0.06)); label("$"+string(i)+"$",(i,0),S,fontsize(7pt)); draw((-0.06,i)--(0.06,i)); label("$"+string(i)+"$",(0,i),W,fontsize(7pt));}
draw((0,0)--(L-0.2,L-0.2),gray+linewidth(0.8));
pair f(real k){real x=k^(1/(k-1)); return (x,k*x);}
guide g1, g2;
for(int i=0;i<=200;++i){real k=1.0005+i*0.02; pair p=f(k); if(p.y<=L-0.1){g1=g1--p; g2=g2--(p.y,p.x);}}
draw(g1,blue+linewidth(0.9)); draw(g2,blue+linewidth(0.9));
dot((exp(1),exp(1)),red); label("$(e,e)$",(4.7,3.3),red+fontsize(8pt)); draw((4.3,3.25)--(2.85,2.75),red,Arrow(4));
dot((2,4)); dot((4,2)); label("$(2,4)$",(2,4),E,fontsize(7pt)); label("$(4,2)$",(4,2),N,fontsize(7pt));
draw((1,0)--(1,L),dashed+gray); draw((0,1)--(L,1),dashed+gray);
\end{asy}

The sketch shows the line $y=x$ (grey) and the second curve (blue), which passes through $(2,4)$ and $(4,2)$, has the asymptotes $x=1$ and $y=1$, and crosses the line at $(e,e)$.
